\section{Complete quotient moments and uniform transfer}
\label{sec:general-transfer}
For finite $X$ and a complete source $J$ define
\[
 Z_J(X)=\sum_{N\lhd X}|\Epi(J,X/N)|,\qquad
 h(w)=w-(w\bmod2).
\]
The normal subgroups here are literal, including $N=X$. Distinct axes
with isomorphic quotient groups remain distinct terms.

\begin{lemma}[The complete quotient moment]\label{lem:quotient-moment}
If $R$ has a specified faithful permutation representation of degree
$v$, then for every integer $q\geq1$,
\[
 \sum_{J\leq S_b}Z_J(R)^q\leq s_{b+qv}.
\]
\end{lemma}
\begin{proof}
$Z_J(R)$ counts full subdirect subgroups of $R\times J$ by their
intersection with $R$ and their literal graph map. For an ordered
$q$-tuple of such subgroups take their fibre product over the same
$J$ inside $R^q\times J$. Projection recovers $J$ and every member
of the tuple: each of the other subgroups can lift the same element
of $J$. This is an injection on actual counting sets, including
coincident maps and proper joint quotient images.
\end{proof}

\begin{lemma}[Uniform normalization]\label{lem:uniform-delta}
Uniformly for $n=b+w$, $w\geq1$,
\begin{equation}\label{eq:uniform-delta}
 \log\Delta(n,b)\leq-\frac{h(w)b}{8}-\frac{w^2}{16}
                                      +O(w\log(n+2)).
\end{equation}
\end{lemma}
\begin{proof}
Let $r=\lfloor b/2\rfloor$ and $R=\lfloor n/2\rfloor$.
The recurrence gives $c_r/c_R\leq(2R)^{R-r}$, and
$p_b\leq(1+b/6)c_r$, whereas $p_n\geq c_R$.
These terms cost $O(w\log(n+2))$, since $R-r\leq(w+1)/2$.
The Gaussian bounds give
$\log(G_r/G_R)=(r^2-R^2)/4+O(1)$.
For even $w$ this is $-wb/8-w^2/16+O(w)$.
For odd $w$, the parity correction is at most $b/8+O(w)$,
giving \eqref{eq:uniform-delta}. The smallest indices are absorbed
in the same absolute constant.
\end{proof}

\begin{theorem}[A finite complete comparator]\label{thm:finite-comparator}
Fix a finite family of actual transitive actions $U$ of degrees $w$.
Suppose each has a fixed faithfully represented comparator $R_U$ of
degree $0<v<h(w)$ and constants $B>0$, $C\geq0$, $\eta<h(w)/8$
such that for all complete $J\leq S_b$,
\[
 Z_J(U)\leq BZ_J(R_U)+C2^{\eta b}.
\]
Then their complete pointed union has one quadratic-exponentially
small scalar and one exponentially small complete forward row, with
the original action weights. The comparator need not be a quotient
of $U$.
\end{theorem}
\begin{proof}
Put $\alpha=(v+h)/16$, $e=(h-v)/16$ and
$\chi=h/8-\max(\alpha,\eta)>0$. Split at
$Z_J(R)>2^{\alpha b}$. The cold summand is exactly bounded by
\[
 \frac{\Delta(n,b)}{a(U)}
           (B2^{\alpha b}+C2^{\eta b}).
\]
Equation~\eqref{eq:uniform-delta} gives exponential row decay after
the finite sum. On the hot part put $d=\max(\eta-\alpha,0)$, so
the whole fibre is at most $(B+C)2^{db}Z_J(R)$.
Use Lemma~\ref{lem:quotient-moment} at
$q=\max(1,\lceil8eb/v^2\rceil)$, $q_*=8eb/v^2$.
The main normalized exponent is
\[
 \frac{(b+qv)^2}{16}-(q-1)\alpha b+db-\frac{(b+w)^2}{16}
 =-\frac{4e^2b^2}{v^2}-\chi b-\frac{(w-h)b}{8}
   -\frac{w^2}{16}+\frac{v^2}{16}(q-q_*)^2.
\]
Here $0\leq q-q_*\leq1$ and all parameters are fixed. The coarse
error in \eqref{eq:coarse} and the original factorials cost
$O(n^{3/2}+n\log(n+2))$, smaller than the quadratic loss.
Take one any-hot union and its cold complement as in
Theorem~\ref{thm:fusion}.
\end{proof}

\begin{theorem}[Growing-parameter comparator transfer]\label{thm:parametric}
Fix $w_0\geq3$ and $0<\rho\leq1/8$.
For each $w\geq w_0$, let $\Sigma_w$ be a finite certificate menu.
Each $\sigma\in\Sigma_w$ carries an actual transitive action $U_\sigma$
of degree $w$, a finite comparator $R_\sigma$, a faithful action of
$R_\sigma$ on $v_{0,\sigma}$ points, a relabelling-invariant acceptance
predicate on its complete physical fibres, and $\eta_\sigma\geq0$.
The same action may occur with several certificate labels. Require
\[
 \frac{h(w)-v_{0,\sigma}}{8}-\eta_\sigma\geq\rho w.
\]
Write $Z_{J,\sigma}(U_\sigma)$ for the literal quotient-map weight of
the accepted part of the complete source $J$. Suppose that, for every
$J\leq S_b$,
\begin{equation}\label{eq:parametric-fibre}
 Z_{J,\sigma}(U_\sigma)\leq
 D_\sigma(b)2^{\eta_\sigma b}Z_J(R_\sigma),
 \qquad D_\sigma(b)\geq0.
\end{equation}
Write $a(U_\sigma)$ for the original action normalizer and assume the direct
original-weight menu bound
\begin{equation}\label{eq:parametric-menu-mass}
 \text{for every $\gamma>0$, eventually uniformly for $w_0\leq w\leq n$,}
 \qquad
 \sum_{\sigma\in\Sigma_w}
       \frac{D_\sigma(n-w)}{a(U_\sigma)}\leq2^{\gamma wn}.
\end{equation}
Then the complete accepted physical union has a forward kernel and scalar
with the uniform eventual bounds
\begin{equation}\label{eq:parametric-rates}
 K(n,n-w)\leq2^{-\rho wn/8},\qquad
 \sum_bK(n,b)\leq2\cdot2^{-\rho w_0n/8},\qquad
 \eta(n)\leq2^{-\rho^2n^2/8}.
\end{equation}
\end{theorem}
\begin{proof}
Pad the given representation by fixed points to
$v=\max(v_{0,\sigma},\lfloor4\rho h(w)\rfloor)$ and put
\[
 \delta=(h(w)-v)/8-\eta_\sigma,\qquad a=v/8+\delta/2.
\]
Then $\delta\geq\rho w/2$, $v\leq(1-4\rho)w$, and
$\eta_\sigma+a=h(w)/8-\delta/2$. If $R_\sigma\ne1$, also
$v\geq\rho w$: indeed the original margin and $\eta_\sigma\geq0$ give
$v_{0,\sigma}\leq h(w)-8\rho w\leq(1-4\rho)w$, while
$\lfloor4\rho h(w)\rfloor\leq(1-4\rho)w$.
The loss from padding is at most $4\rho w/8$, proving the stated lower
bound on $\delta$. Faithfulness gives $v_{0,\sigma}\geq2$ when
$R_\sigma\ne1$; use this
when $\rho w\leq1$, and otherwise
$\lfloor4\rho h(w)\rfloor\geq4\rho w-4\rho-1>\rho w$.
For $R_\sigma=1$ the hot event below is empty, so no lower bound on $v$ is
needed. This controlled padding preserves the original margin.

Split at $Z_J(R_\sigma)>2^{ab}$. The cold kernel is
\[
 K(n,b)=\Delta(n,b)\sum_{\sigma\in\Sigma_{n-b}}
         \frac{D_\sigma(b)2^{(\eta_\sigma+a)b}}{a(U_\sigma)}.
\]
By \eqref{eq:uniform-delta} its main exponent is at most
$-\rho wb/4-w^2/16\leq-\rho wn/4-w^2/32$.
The estimate is uniform from $w=3$: if $r=\lfloor w/2\rfloor$, then
$w\leq3r$ and hence $3w^2\leq32r^2$; together with
$\rho\leq1/8$ this supplies the floor-term comparison used in the last
inequality.
The number of transitive action classes is $2^{o(w^2)}$ by the
transitive-subgroup bound of \cite{LucchiniMenegazzoMorigi2000}, in
the labelled form quoted in the introduction of
\cite{RoneyDougalTracey2025}. Thus the formerly used sufficient
conditions
\[
 D_\sigma(b)\leq2^{\beta(w)+Cw\log^2(w+b+1)},
 \qquad \beta(w)=o(w^2),\qquad
 \max_U|\{\sigma:U_\sigma=U\}|=2^{o(w^2)}
\]
imply \eqref{eq:parametric-menu-mass}: the logarithms of the class count
and the certificate multiplicity, together with $\beta(w)$, are $o(w^2)$,
while the logarithmic term is $o(wn)$ uniformly for
$w_0\leq w\leq n$. In the present formulation the direct menu bound
absorbs these terms together with the pointing polynomial. This proves
the first bound in \eqref{eq:parametric-rates}. Sum the geometric tail
over $w\geq w_0$ to obtain the second.

For a nonempty hot part use
$q=\max(1,\lceil4\delta b/v^2\rceil)$, $q_*=4\delta b/v^2$,
and $M=b+qv$. Lemma~\ref{lem:quotient-moment} gives
$\sum_{J\ {\rm hot}}Z_J(R_\sigma)\leq2^{-(q-1)ab}s_M$.
The exact main exponent, including the lift factor, is
\begin{align*}
 \frac{M^2}{16}-(q-1)ab+\eta_\sigma b-\frac{n^2}{16}
 &=-\frac{\delta^2b^2}{v^2}
   +(\eta_\sigma+a-w/8)b-\frac{w^2}{16}
                      +\frac{v^2}{16}(q-q_*)^2\\
 &\leq-\frac{\rho^2b^2}{4}-\frac{\rho wb}{4}
                 -(\rho/2-\rho^2)w^2
 \leq-\rho^2n^2/4.
\end{align*}
The last comparison uses $\rho\leq1/8$; the adverse odd correction
$(h(w)-w)b/8$ is nonpositive. Moreover
$M\leq(1+1/(2\rho))b+w$, so the coarse error is uniformly
$O_\rho(n^{3/2})$. The original hot sum is
\[
 \sum_{w=w_0}^n\sum_{\sigma\in\Sigma_w}
 \frac{D_\sigma(b)2^{\eta_\sigma b}}
      {b!a(U_\sigma)p_nG_{\lfloor n/2\rfloor}}
       \sum_{J\ {\rm hot}}Z_J(R_\sigma),\qquad b=n-w.
\]
Its remaining normalization costs $O(n\log(n+2))$.
Equation~\eqref{eq:parametric-menu-mass}, first at each width and then
after the sum over at most $n+1$ widths, absorbs the original weighted
menu into an arbitrarily small multiple of $n^2$. Retaining half the
quadratic reserve proves the last bound in
\eqref{eq:parametric-rates}. No normalized-count boundedness has been
assumed.
\end{proof}

\begin{lemma}[Concrete menu-mass certificate]
\label{lem:concrete-menu-mass}
Suppose $w_0\geq1$ and that, for every $\epsilon>0$, there are
$B_\epsilon\geq0$ and $p_\epsilon\in\mathbb N$ such that, uniformly for
$w_0\leq w\leq n$,
\[
 \sum_{\sigma\in\Sigma_w}
       \frac{D_\sigma(n-w)}{a(U_\sigma)}
 \leq B_\epsilon(n+1)^{p_\epsilon}2^{\epsilon w^2}.
\]
Then the direct menu bound \eqref{eq:parametric-menu-mass} holds.
\end{lemma}
\begin{proof}
Given $\gamma>0$, take $\epsilon=\gamma/3$.  Eventually the fixed constant
and the fixed polynomial are each at most $2^{\epsilon n}$.  Since
$1\leq w\leq n$,
\[
 2\epsilon n+\epsilon w^2\leq\gamma wn.
\]
The threshold depends on $\gamma$, but is uniform in $w$ on the indicated
triangle.
\end{proof}

\begin{lemma}[Index--entry menu aggregation]
\label{lem:index-entry-menu-mass}
Suppose that, for every $\epsilon>0$, there are constants
$B_{1,\epsilon},B_{2,\epsilon}\geq0$ and
$p_\epsilon\in\mathbb N$ such that
\[
 |\Sigma_w|\leq B_{1,\epsilon}2^{\epsilon w^2}
\]
and, uniformly for $w_0\leq w\leq n$ and $\sigma\in\Sigma_w$,
\[
 \frac{D_\sigma(n-w)}{a(U_\sigma)}
 \leq B_{2,\epsilon}(n+1)^{p_\epsilon}2^{\epsilon w^2}.
\]
Then the hypothesis of Lemma~\ref{lem:concrete-menu-mass} holds.
\end{lemma}
\begin{proof}
Apply the two hypotheses with $\epsilon/2$ and sum the common entry bound
over $\Sigma_w$.  The constants multiply, the polynomial degree is unchanged,
and the two width exponents add to $\epsilon w^2$.
\end{proof}

\begin{lemma}[Retained-cell incidence]\label{lem:retained-cell-incidence}
For a complete source $J$ and finite comparator $R$, put
\[
 \mathcal Q_J(R)=\coprod_{N\lhd R}\Epi(J,R/N),
 \qquad |\mathcal Q_J(R)|=Z_J(R).
\]
Let $\mathcal E$ be an accepted intrinsic carrier-epimorphism family and
let $A$ be a finite set of retained invariant flags.  If
\[
 \kappa:\mathcal E\longrightarrow\mathcal Q_J(R)\times A
 \quad\hbox{has every fibre of size at most }L,
\]
then
\begin{equation}\label{eq:retained-cell-incidence}
 |\mathcal E|\leq |A|LZ_J(R).
\end{equation}
Consequently $|A|L\leq D(b)2^{\eta b}$ proves
\eqref{eq:parametric-fibre}.  The same conclusion holds on the original
literal normal axis whenever the carrier transport reconstructs the original
subgroup and acceptance is tested after that inverse reconstruction.
\end{lemma}
\begin{proof}
Partition $\mathcal E$ into the fibres of $\kappa$ and sum the bound $L$
over the $Z_J(R)|A|$ possible cells.  Reversible transport is injective on
the original axis, preserves the complete exterior source $J$, and tests the
same original acceptance predicate, so the transported family is a subset
of $\mathcal E$ with no additional multiplicity.

For the capacity input attached to an epimorphism
$\delta:J\twoheadrightarrow Q$, use the literal pullback
\[
 H_\delta=\{(u,j):\beta(u)=\delta(j)\}\leq U\times J.
\]
Both projections are full and its first axis is exactly $\ker\beta$.
Consequently the joint character, derived-normal, invariant-radical and
restricted-transgression inequalities apply to this same $H_\delta$, even
when it is a nonabelian proper subdirect product.  No direct-product
replacement or separately maximized axis is introduced before forming the
retained cell.

The subgroup $H_\delta$ also recovers $\delta$: surjectivity of $\beta$
lets equality of two pullbacks be tested above every value of the common
quotient.  Hence, for any preliminary comparator map $\kappa_0$, the exact
cell
\[
 \delta\longmapsto\bigl(\kappa_0(\delta),H_\delta\bigr)
\]
is injective and has fibres of size at most one.  A radical or restricted
annihilator flag may replace $H_\delta$ only together with a proved bound for
the fibres created by that coarsening; this bound is precisely the $L$ in
\eqref{eq:retained-cell-incidence}.
\end{proof}

The repeated action labels in Theorem~\ref{thm:parametric} implement the
owner-or-capacity dichotomy.  Take a finite ordered list of predicates
invariant under relabelling of the complete physical subgroup and assign each
subgroup to its first eligible predicate.  Pair that owner with the single
family orbit carrying its certificate whenever such an orbit exists; only
the final catch-all branch uses an arbitrary selected bad orbit.  First
ownership makes these pairs disjoint, while the local action check makes an
owner paired with an unrelated orbit an empty cell.  Hence the catch-all
branch holds precisely when no family orbit exists.  Earlier owners and the
retained-annihilator branch may use different local certificates
$\sigma=(\text{owner},U)$; only their combined original-weight mass in
\eqref{eq:parametric-menu-mass} is charged.  The
complete one-triple family is the first required application of these same
local bounds.  In its independent finite-owner audit, full projection onto
the displayed outer $C_3$ and the retained one-$C_3$ source pattern remove the
apparent degree-three ternary-owner label before the selected carrier is
forgotten into its canonical cell.

\begin{proposition}[Pre-$E_7$ owner-or-cell transfer]
\label{prop:pre-e7-owner-cell-transfer}
Restrict the first-owner action menu to the nonbinary orbit actions that admit
no selected $E_7$ pair frame.  On a literal normal axis $N$, suppose either
an earlier owner bounds the surviving epimorphisms by the common complete
quotient-map weight (as a quotient comparator does), or there is a reversible
carrier and a retained-cell map as in
Lemma~\ref{lem:retained-cell-incidence}, with flag set $A_N$, cell-fibre bound
$L_N$, and
\[
 |A_N|L_N\leq C_N(b)2^{\eta b}.
\]
Give an owned axis the extracted coefficient $2^{-\eta b}$ and an unowned
axis the coefficient $C_N(b)$.  Suppose each resulting complete normal-axis
sum, divided by its original action normalizer, satisfies the uniform
one-entry hypothesis of Lemma~\ref{lem:index-entry-menu-mass}, and the
aggregation parameters satisfy the hypotheses of
Theorem~\ref{thm:parametric}.  Then the exact pre-$E_7$ physical family has
the forward kernel and scalar bounds \eqref{eq:parametric-rates}.
\end{proposition}
\begin{proof}
On an owned axis the direct owner envelope bounds the surviving epimorphisms
by the complete quotient-map weight.  Multiplication by the extracted common
factor $2^{\eta b}$ cancels $2^{-\eta b}$, so this contribution has coefficient
exactly one.  On every other axis reversibility returns the retained-cell
bound to the same original source and normal axis, and
Lemma~\ref{lem:retained-cell-incidence} gives coefficient
$C_N(b)2^{\eta b}$.  Sum over the literal normal axes before dividing by the
original action normalizer.  The transitive-subgroup theorem gives
$2^{o(w^2)}$ labelled actions.  Passing to conjugacy classes and imposing the
pre-$E_7$ restriction can only decrease their number, while the fixed
first-owner menu adds only a constant factor.  Lemma~\ref{lem:index-entry-menu-mass}
therefore supplies Lemma~\ref{lem:concrete-menu-mass}, which gives
\eqref{eq:parametric-menu-mass}.  Theorem~\ref{thm:parametric} now applies to
the restricted physical cover.
\end{proof}

The owner order in this proposition serves only to make the physical
partition disjoint.  None of the local estimates needs the assertion that an
earlier predicate failed: one may choose any valid certificate for the
retained action class.  Formally, the family predicate says that such a
certificate exists on a literal orbit.  Transporting that orbit under an
ambient relabelling proves degree naturality, and the cover displays it before
forming the local fusion cell.  The residual catch-all may therefore infer
the absence of every family orbit without a separate precedence theorem.

The following additive form is used for families whose sharp estimate has a
source-independent remainder.  It avoids forcing that remainder through the
complete-quotient threshold.

\begin{proposition}[Additive-tail pre-$E_7$ transfer]
\label{prop:pre-e7-additive-tail-transfer}
In the setting of Proposition~\ref{prop:pre-e7-owner-cell-transfer}, suppose
the literal normal axis $N$ has the unconditional envelope
\[
 Z_{J,N}\leq C_N(b)2^{\eta b}Z_J(R)
                 +T_N(b)2^{\theta b},
 \qquad C_N(b),T_N(b)\geq0.
\]
Sum $C_N$ and $T_N$ over the literal normal axes before applying the original
action-normalizer divisor.  Assume both resulting weighted menus satisfy
\eqref{eq:parametric-menu-mass}, the main parameters satisfy
Theorem~\ref{thm:parametric}, and uniformly
\[
 \theta\leq \frac{h(w)}8-\frac{\rho w}4.
\]
Then this family has a complete exponentially contractive forward estimate.
The scalar is exactly the scalar of the comparator-weighted main term; the
additive remainder contributes only a cold row.
\end{proposition}
\begin{proof}
For fixed $J$, sum the displayed inequality over the original normal axes.
Apply the complete quotient moment and its hot/cold split only to the first
sum.  The second sum is independent of $J$, so summing it over all complete
$J\leq S_b$ gives
\[
 s_b\Bigl(\sum_NT_N(b)\Bigr)2^{\theta b}.
\]
After the exact pointing and benchmark normalization this is the usual cold
kernel with exponent $\theta$.  Equation~\eqref{eq:uniform-delta}, the
displayed gap, and the second menu-mass hypothesis give an exponentially
decaying row.  Add this cold-only forward estimate to the estimate supplied
by Theorem~\ref{thm:parametric}.  No quotient threshold, hot scalar, or
independent worst-case value of $Z_J(R)$ is introduced for the tail.
\end{proof}

\begin{lemma}[Central nonbinary pair sections]\label{lem:central-pair}
For a faithful nonbinary transitive action of even degree $s\geq4$,
every trivial section of its binary permutation module has dimension
$t\leq s/2-1$. Consequently a central pair-kernel quotient at original
width $2s$ has an effective graph moment degree $s+2t\leq2s-2$.
\end{lemma}
\begin{proof}
If an odd prime $p$ divides $s$, its Sylow subgroup has no fixed
points. Exactness of invariants bounds a trivial section by the number
of its orbits, at most $s/p\leq s/2-1$.
If $s=2^a\geq8$, a Sylow $2$-subgroup is transitive. The soluble
permutation-module bound gives $t\leq B_a\leq3s/8\leq s/2-1$.
At $s=4$, the $3$-cycle fixed-vector argument in
Lemma~\ref{lem:unipotent} gives $t\leq1$.
For a central kernel $A$ take $C=A$ in Proposition~\ref{prop:prefix}.
Then the cocycle tail vanishes and the moment degree is $s+2t$.
\end{proof}
